\documentclass[../main.tex]{subfiles}
\begin{document}

\iniciarSecao[Resultados HIL]{Metodologia Experimental e Resultados HIL}

% ---------------------------------------------------------------------------
% Matriz ADR-020 e Resultados Empiricos
% ---------------------------------------------------------------------------
\begin{frame}
  \frametitle{Matriz Experimental ADR-020 e Resultados HIL (1024 B, $N=20$)}
  \scriptsize
  \renewcommand{\arraystretch}{1.15}
  \begin{tabular}{@{}llccccc@{}}
    \toprule
    \textbf{Cen\'ario} & \textbf{Perfil} & \textbf{Lat\^encia (ms)} & \textbf{Ciclos PMU} & \textbf{Ciclos PL} & \textbf{CR (\%)} & \textbf{Speedup}\\
    \midrule
    1: SW Puro & Quiescent & 3,796 & 2.467.688 & 0 & +55,7\% & 1,00x\\
    1: SW Puro & Dynamic & 10,188 & 6.621.969 & 0 & +2,5\% & 1,00x\\
    1: SW Puro & Chaotic & 12,907 & 8.389.708 & 0 & -12,6\% & 1,00x\\
    \midrule
    2: PL LZSS + SW AES & Quiescent & 5,990 & 3.893.369 & 200.972 & +55,7\% & 0,63x\\
    2: PL LZSS + SW AES & Dynamic & 16,112 & 10.473.009 & 573.622 & +2,5\% & 0,63x\\
    2: PL LZSS + SW AES & Chaotic & 20,563 & 13.366.113 & 744.661 & -12,6\% & 0,63x\\
    \midrule
    \textbf{3: SW LZSS + PL AES} & \textbf{Quiescent} & \textbf{3,068} & \textbf{1.994.017} & \textbf{668} & +55,7\% & \textbf{1,24x}\\
    \textbf{3: SW LZSS + PL AES} & \textbf{Dynamic} & \textbf{8,579} & \textbf{5.576.620} & \textbf{1.452} & +2,5\% & \textbf{1,19x}\\
    \textbf{3: SW LZSS + PL AES} & \textbf{Chaotic} & \textbf{11,039} & \textbf{7.175.201} & \textbf{1.679} & -12,6\% & \textbf{1,17x}\\
    \midrule
    4: HW Total & Quiescent & 5,260 & 3.419.303 & 201.639 & +55,7\% & 0,72x\\
    4: HW Total & Dynamic & 14,503 & 9.427.015 & 575.075 & +2,5\% & 0,70x\\
    4: HW Total & Chaotic & 18,694 & 12.150.877 & 746.340 & -12,6\% & 0,69x\\
    \bottomrule
  \end{tabular}
  \fonte{Matriz\_de\_Testes\_e\_M\'etricas.csv, README.md, ADR-020}
\end{frame}

% ---------------------------------------------------------------------------
% Recursos FPGA e Achados (Risco R9)
% ---------------------------------------------------------------------------
\begin{frame}
  \frametitle{Recursos da FPGA XC7Z010 e Confirma\c{c}\~ao do Risco R9}
  \begin{columns}[T]
    \begin{column}{.48\textwidth}
      \begin{block}{Ocupa\c{c}\~ao P\'os-Implementa\c{c}\~ao (Cen\'ario 4)}
        \begin{itemize}
          \item \textbf{LUTs:} 8.810 de 17.600 (50,06\%).
          \item \textbf{Flip-Flops:} 12.048 de 35.200 (34,23\%).
          \item \textbf{Block RAM:} 9 de 60 RAMB36 (15,00\%).
          \item \textbf{DSP48E1:} 0 de 80 (0,00\%).
          \item \textbf{Timing:} \texttt{WNS = +0,551 ns}, \texttt{WHS = +0,008 ns}, 0 latches.
        \end{itemize}
      \end{block}
    \end{column}
    \begin{column}{.48\textwidth}
      \begin{pontoschave}[Achados Emp\'iricos]
        \begin{itemize}
          \item \textbf{Risco R9 confirmado:} busca sequencial na janela LZSS via Porta B a 100 MHz (201k a 745k ciclos) mais c\'opia AXI/BRAM torna o PL LZSS mais lento (0,63x) que o Cortex-A9 com cache L1/L2 a 650 MHz.
          \item \textbf{Ganho do Cen\'ario 3:} \texttt{SW LZSS + PL AES-CTR} obt\'em \textbf{1,17x a 1,24x} de speedup ponta a ponta em todos os perfis.
          \item Paridade bit-a-bit validada nos 4 cen\'arios.
        \end{itemize}
      \end{pontoschave}
    \end{column}
  \end{columns}
  \fonte{README.md, STATUS.md, Matriz\_de\_Testes\_e\_M\'etricas.csv}
\end{frame}

\end{document}
